\section{Экономическое обоснование разработки и реализации персонального ассистента организации времени на массовом рынке}
\label{sec:economics}

\subsection{Характеристика персонального ассистента организации времени, разрабатываемого для реализации на рынке}


Целью дипломного проекта является разработка персонального ассистента организации времени, предназначенного для поддержки обучающихся при планировании учебной деятельности и распределении нагрузки по календарю.

Разрабатываемое программное средство представляет собой прикладной программный продукт, ориентированный на использование в образовательной сфере. Ассистент обеспечивает естественно-языковое взаимодействие с пользователем, автоматическую декомпозицию учебных задач и формирование календарного плана с учетом занятых временных интервалов.

Область применения продукта включает:
\begin{itemize}
    \item персональное планирование учебной нагрузки;
    \item организацию самостоятельной работы;
    \item контроль сроков выполнения заданий;
    \item формирование индивидуального графика учебной активности.
\end{itemize}

Основную целевую аудиторию продукта, разрабатываемого <<ООО Фабрика Инноваций и Решений>>, составляют обучающиеся учреждений высшего образования. Потенциальными покупателями также могут выступать образовательные организации, заинтересованные во внедрении цифровых сервисов поддержки учебного процесса.

К основным функциям программного продукта относятся:
\begin{itemize}
    \item прием пользовательского запроса через интерфейс взаимодействия;
    \item анализ содержания поставленной цели;
    \item формирование набора подзадач;
    \item автоматическое распределение задач по календарю;
    \item отображение чернового варианта расписания;
    \item ручная корректировка предложенного плана;
    \item фиксация утвержденного результата.
\end{itemize}

В отличие от обычных календарей и списков задач персональный ассистент организации времени ориентирован на академический контекст и сочетает автоматизацию планирования с возможностью ручной корректировки результата. Это снижает когнитивную нагрузку на пользователя и делает процесс планирования более последовательным.

Разрабатываемый продукт целесообразно рассматривать как товар, реализуемый на массовом рынке информационных технологий, поскольку он не создается под требования одного заказчика и может тиражироваться для различных категорий пользователей без изменения базовой архитектуры.

Возможные модели коммерциализации включают:
\begin{itemize}
    \item индивидуальную подписку для конечных пользователей;
    \item институциональный доступ для университетов и иных образовательных организаций;
    \item лицензирование продукта для организованных групп обучающихся.
\end{itemize}

Актуальность продукта определяется распространенностью проблем самоорганизации обучающихся, необходимостью управления дедлайнами и спросом на цифровые инструменты личной продуктивности. На рынке присутствуют календарные сервисы, менеджеры задач и системы интеллектуального планирования, однако большинство таких решений не ориентировано специально на учебный процесс, что приводится в аналитическом разделе. Это позволяет рассматривать персональный ассистент организации времени как специализированный программный продукт, обладающий рыночным потенциалом.

При экономическом обосновании проекта персональный ассистент организации времени рассматривается как программный продукт с продуктовой бизнес-моделью. Предполагается, что его распространение будет осуществляться через цифровые каналы доступа, а экономический результат для организации-разработчика будет выражаться в приросте чистой прибыли от реализации лицензий или подписок.


\subsection{Расчет инвестиций в разработку программного средства для его реализации на рынке}

Инвестициями в разработку персонального ассистента организации времени являются затраты, необходимые для создания и подготовки программного продукта к выводу на рынок. В состав указанных затрат целесообразно включить основную заработную плату разработчиков, дополнительную заработную плату, отчисления на социальные нужды, прочие расходы, а также расходы на реализацию программного продукта. Такой состав затрат соответствует методике экономического обоснования программного продукта, предназначенного для свободной реализации на рынке информационных технологий \cite{Economics}.

На первом этапе необходимо определить затраты на основную заработную плату команды разработчиков. Расчет выполняется исходя из состава команды, уровня месячной заработной платы исполнителей, их часовых ставок и трудоемкости выполняемых работ по проекту, по формуле:

\begin{equation}
    \text{З}_{\text{о}}=\text{К}_{\text{пр}}\cdot \sum_{i=1}^{n}{\text{З}_{\text{ч}i}\cdot t_{i}},
\end{equation}

где $\text{К}_{\text{пр}}$ -- коэффициент премий (равный 1.5);

    $n$ -- категории исполнителей, занятых разработкой программного средства;

    $\text{З}_{\text{ч}i}$ -- часовая заработная плата исполнителя i-й категории, р.;

    $t_{i}$ -- трудоемкость работ, выполняемых исполнителем i-й категории, определяется исходя из сложности разработки программного обеспечения и объема выполняемых им функций, ч.

% После представления формулы необходимо привести расшифровку используемых обозначений и указать принятый состав команды разработки.

Для персонального ассистента организации времени в состав исполнителей потребуется включить следующих специалистов: архитектор или ведущий разработчик, инженер-программист, специалист по тестированию, дизайнер интерфейсов и аналитик. Значения заработной платы приняты на основе открытых обзоров рынка труда и агрегаторов зарплат по ИТ-специальностям в Минске и Беларуси \cite{RecruitmentBy2026,LevelsFYI2026,WorldSalariesAnalyst2026,WorldSalariesQA2026,WorldSalariesUX2026}.

% Конкретный набор ролей и значения трудоемкости определяются на основании фактического содержания работ по проекту.

\begin{table}[!h!t]
\caption{Расчет затрат на основную заработную плату разработчиков}
\label{tab1}
\centering
    \begin{tabular}{| >{\raggedright\arraybackslash}m{0.2\textwidth}
            | >{\centering\arraybackslash}m{0.15\textwidth}
            | >{\centering\arraybackslash}m{0.15\textwidth}
            | >{\centering\arraybackslash}m{0.20\textwidth}
            | >{\centering\arraybackslash}m{0.15\textwidth}|}
\hline
Категория исполнителя & Месячный оклад, р. & Часовой оклад, р. & Трудоемкость работ, ч & Итого, р. \\
\hline
Архитектор решений & 8200.00 & 48.80 & 120 & 8784.00 \\ \hline
Бизнес-аналитик & 4800.00 & 28.57 & 168 & 7199.64 \\ \hline
Дизайнер & 3200.00 & 19.05 & 120 & 3429.00 \\ \hline
Инженер-программист & 7100.00 & 42.26 & 720 & 45640.80 \\ \hline
Тестировщик & 4000.00 & 23.80 & 230 & 8211.00 \\ \hline
\multicolumn{4}{|l|}{Всего затрат на основную заработную плату разработчиков} & 73264.44 \\
\hline
\end{tabular}
\end{table}

Следующим этапом рассчитывается дополнительная заработная плата разработчиков, представляющая собой выплаты, предусмотренные законодательством и внутренними нормативами, в процентах к основной заработной плате.

\begin{equation}
    \text{З}_{\text{д}} = \frac{\text{З}_{o}\cdot H_\text{д}}{100},
\end{equation}

где $\text{З}_{\text{o}}$ -- затраты на основную заработную плату, р.;

    $H_\text{д}$ -- норматив дополнительной заработной платы (20 \%).

Подставим значение в формулу и вычислим $\text{З}_{\text{д}}$:
$$
\text{З}_{\text{д}} = \frac{73264.44 \cdot 20}{100} = 14652.88 \text{ р}.
$$

Далее определяются отчисления на социальные нужды, включающие обязательные начисления на фонд оплаты труда в соответствии с действующим законодательством.

\begin{equation}
    \text{P}_{\text{соц}} = \frac{(\text{З}_{o} + \text{З}_{\text{д}})\cdot H_\text{соц}}{100},
\end{equation}

где $H_\text{соц}$ -- норматив отчислений на социальные нужды, (34.6 \%).

Подставим результаты вычислений в формулу и вычислим $\text{P}_{\text{соц}}$:
$$
 \text{P}_{\text{соц}} = \frac{(73264.44 + 14652.88)\cdot 34.6}{100} = 30419.39 \text{ р}.
$$
После этого рассчитываются прочие расходы, связанные с обеспечением процесса разработки. К ним могут относиться затраты на организацию рабочего процесса, использование оборудования и инфраструктуры, административное сопровождение и иные косвенные расходы, принимаемые по нормативу в процентах от основной заработной платы.

\begin{equation}
    \text{P}_{\text{пр}} = \frac{\text{З}_{o} \cdot \text{H}_{\text{пз}}}{100},
\end{equation}

где $\text{H}_{\text{пз}}$ -- норматив прочих затрат в целом по организации, (30 \%).

Подставим значение из выражения в формулу и произведем расчет $\text{P}_{\text{пр}}$:
$$
 \text{P}_{\text{пр}} = \frac{73264.44 \cdot 30}{100} = 21979.33 \text{ р}.
$$

Для программного продукта, ориентированного на массовый рынок, дополнительно учитываются расходы на реализацию. В эту статью могут входить затраты, связанные с выводом продукта на рынок, обеспечением его распространения и сопровождением каналов продаж.

\begin{equation}
    \text{P}_{\text{р}} = \frac{\text{З}_{o} \cdot \text{H}_{\text{р}}}{100},
\end{equation}

где $\text{H}_{\text{р}}$ -- норматив расходов на реализацию, (4 \%).

Подставим значение из выражения в формулу и произведем расчет $\text{P}_{\text{р}}$:
$$
 \text{P}_{\text{р}} = \frac{73264.44 \cdot 4}{100} = 2930.57 \text{ р}.
$$

На основании рассчитанных статей определяется общая сумма затрат на разработку и реализацию программного продукта.


\begin{equation}
    \text{З}_{\text{р}} = \text{З}_{o} + \text{З}_{\text{д}} + \text{P}_{\text{соц}} + \text{P}_{\text{пр}} + \text{P}_{\text{р}}
\end{equation}

Подставим результаты вычислений в формулу и произведем расчет:
$$
    \text{З}_{\text{р}} = 73264.44 + 14652.88 + 30419.39 + 21979.33 + 2930.57 = 143246.61 \text{ р}.
$$


Итоговые результаты расчета инвестиций целесообразно представить в сводной таблице, включающей все перечисленные статьи затрат.


\begin{table}[!h!t]
\centering
\caption{Результаты расчета цены на разработку программного средства}
\label{tab2}
\begin{tabular}{|l|c|}
\hline
\multicolumn{1}{|c|}{Наименование статьи затрат} & Сумма, р. \\ \hline
1 Основная заработная плата разработчиков & 73264.44 \\ \hline
2 Дополнительная заработная плата разработчиков & 14652.88 \\ \hline
3 Отчисления на социальные нужды & 30419.39 \\ \hline
4 Прочие расходы & 21979.33 \\ \hline
5 Расходы на реализацию & 2930.57 \\ \hline
6 Общая сумма затрат на 
разработку и реализацию & 143246.61  \\ \hline
\end{tabular}
\end{table}


Таким образом, в данном подразделе должна быть получена совокупная величина инвестиций, необходимая для разработки персонального ассистента организации времени и его подготовки к рыночной реализации. Полученное значение в дальнейшем используется для сопоставления с прогнозируемым экономическим результатом от продажи программного продукта.

\subsection{Расчет экономического эффекта от реализации программного средства на рынке}

Экономический эффект для организации-разработчика при реализации персонального ассистента организации времени на рынке представляет собой прирост чистой прибыли, получаемой от продажи лицензий или подписок пользователям. Величина этого результата определяется предполагаемым объемом продаж, ценой реализации продукта и принятыми параметрами рентабельности продаж.

Для расчета экономического эффекта необходимо предварительно обосновать предполагаемый годовой объем продаж. При этом следует учитывать специфику целевой аудитории и возможные каналы коммерциализации продукта. В рамках рассматриваемого проекта базовыми сценариями реализации могут выступать индивидуальная подписка для конечных пользователей и предоставление группового доступа образовательным организациям. При выборе прогнозного объема продаж необходимо исходить из консервативной оценки потенциального спроса, чтобы итоговые расчеты сохраняли реалистичность. Дополнительным аргументом в пользу выбранного масштаба реализации является значительная численность обучающихся в системе высшего образования Республики Беларусь, зафиксированная в официальных статистических материалах \cite{BelstatEducation2024}.

Также необходимо определить цену копии, лицензии или подписки на программный продукт. Обоснование цены может осуществляться на основе анализа существующих рыночных аналогов, сравнения доступных моделей монетизации в смежных сервисах цифровой продуктивности и образовательных приложениях, а также с учетом платежеспособности целевой аудитории. Для студенческого сегмента особенно важно выбирать цену, которая будет восприниматься как доступная, а при институциональной модели — учитывать целесообразность закупки продукта университетом для последующего предоставления доступа обучающимся. В качестве ориентиров для сопоставления могут использоваться действующие цены сервисов Todoist, Motion и Sunsama \cite{TodoistPricing2026,MotionPricing2026,SunsamaPricing2026}.

В качестве базового принимается консервативный сценарий реализации программного продукта. Предполагается, что распространение персонального ассистента организации времени будет осуществляться по подписочной модели через веб-платформу. Основным каналом продаж выступают индивидуальные подписки конечных пользователей. Дополнительно допускается приобретение пакетов лицензий образовательными организациями для предоставления доступа обучающимся без необходимости интеграции продукта во внутренние информационные системы. С учетом направленности продукта на студенческую аудиторию, ограниченной платежеспособности конечных пользователей и необходимости обеспечения конкурентоспособности по отношению к существующим цифровым планировщикам годовой объем продаж принимается на уровне 5300 лицензий, что составляет лишь небольшую долю потенциальной аудитории и поэтому может рассматриваться как реалистичная консервативная оценка \cite{BelstatEducation2024}.

Выбранное значение может быть достигнуто за счет сочетания двух каналов продаж: индивидуальных подписок, приобретаемых конечными пользователями, и пакетных лицензий, приобретаемых образовательными организациями в интересах студентов. Такой подход соответствует продуктовой модели распространения и позволяет учитывать как прямой потребительский спрос, так и дополнительный канал реализации без изменения функциональной структуры программного средства.

Отпускная цена одной лицензии принимается равной 200 р. в год, что соответствует примерно 16.67 р. в месяц. Данный уровень цены остается приемлемым для подписочной модели цифрового продукта и может быть обоснован наличием интеллектуальных функций персонального планирования, требующих вычислительных ресурсов и использования больших языковых моделей. При этом выбранная цена остается сопоставимой с ценами зарубежных сервисов цифровой продуктивности и не выходит за пределы умеренного ценового диапазона для годовой подписки \cite{TodoistPricing2026,MotionPricing2026,SunsamaPricing2026}. Таким образом, для дальнейших расчетов принимаются следующие исходные данные: отпускная цена лицензии -- 200 р., годовой объем продаж -- 5300 лицензий, рентабельность продаж -- 35 \%, ставка налога на прибыль -- 18 \%, ставка НДС -- 20 \%.

После обоснования объема продаж и цены программного продукта рассчитывается прирост чистой прибыли, который организация-разработчик получает от его реализации на рынке.

\begin{equation}
    \Delta \text{П}_{\text{ч}}^{\text{р}} = (\text{Ц}_{\text{отп}} \cdot N - \text{НДС}) \cdot \text{Р}_{\text{пр}} \cdot \left(1 - \frac{\text{Н}_{\text{п}}}{100}\right),
\end{equation}

где $\text{Ц}_{\text{отп}}$ -- отпускная цена копии (лицензии) программного средства, р.;

    $N$ -- количество копий (лицензий) программного средства, реализуемое за год, шт.;

    $\text{НДС}$ -- сумма налога на добавленную стоимость, р.;

    $\text{Р}_{\text{пр}}$ -- рентабельность продаж копий (лицензий), доли ед.;

    $\text{Н}_{\text{п}}$ -- ставка налога на прибыль, \%.

Для применения формулы (4.3) предварительно определяется сумма налога на добавленную стоимость, используемая в расчете экономического результата.

\begin{equation}
    \text{НДС} = \frac{\text{Ц}_{\text{отп}} \cdot N \cdot \text{Н}_{\text{д.с}}}{100 + \text{Н}_{\text{д.с}}},
\end{equation}

где $\text{Н}_{\text{д.с}}$ -- ставка налога на добавленную стоимость, \%.

Подставим значения в формулу и определим сумму НДС:
$$
    \text{НДС} = \frac{200.00 \cdot 5300 \cdot 20}{100 + 20} = 176666.67 \text{ р}.
$$

Подставим исходные данные в формулу расчета прироста чистой прибыли:
$$
    \Delta \text{П}_{\text{ч}}^{\text{р}} = (200.00 \cdot 5300 - 176666.67) \cdot 0.35 \cdot \left(1 - \frac{18}{100}\right) = 253516.67 \text{ р}.
$$

Если организация-разработчик является резидентом Парка высоких технологий, методика допускает применение упрощенной формулы расчета прироста чистой прибыли.

\begin{equation}
    \Delta \text{П}_{\text{ч}}^{\text{р}} = \text{Ц}_{\text{отп}} \cdot N \cdot \text{Р}_{\text{пр}}.
\end{equation}

Подставим принятые исходные данные в формулу расчета для резидента ПВТ:
$$
    \Delta \text{П}_{\text{ч}}^{\text{р}} = 200.00 \cdot 5300 \cdot 0.35 = 371000.00 \text{ р}.
$$

Таким образом, при обычном режиме налогообложения прогнозируемый годовой прирост чистой прибыли составляет 253516.67 р., а при условии, что организация-разработчик является резидентом Парка высоких технологий, -- 371000.00 р.

На основании выполненных расчетов в подразделе должна быть получена прогнозная величина годового экономического результата от реализации персонального ассистента организации времени на массовом рынке. Данный показатель является ключевым для последующей оценки экономической эффективности инвестиций, поскольку именно он отражает ожидаемый коммерческий эффект продуктовой модели.

\subsection{Расчет показателей экономической эффективности разработки и реализации программного средства на рынке}

Оценка экономической эффективности разработки и реализации персонального ассистента организации времени зависит от результата сравнения инвестиций в разработку и реализацию продукта с прогнозируемым годовым приростом чистой прибыли.

Поскольку сумма инвестиций в разработку и реализацию программного средства меньше суммы годового прироста чистой прибыли, то есть $143246.61 < 253516.67$, проект окупается менее чем за один год. Следовательно, в соответствии с методикой экономическая эффективность может быть оценена с помощью расчета рентабельности инвестиций.

\begin{equation}
    ROI = \frac{\Delta \text{П}_{\text{ч}}^{\text{р}} - \text{З}_{\text{р}}}{\text{З}_{\text{р}}} \cdot 100 \%,
\end{equation}

где $\Delta \text{П}_{\text{ч}}^{\text{р}}$ -- прирост чистой прибыли, полученной от реализации программного средства на рынке, р.;

    $\text{З}_{\text{р}}$ -- затраты на разработку и реализацию программного средства, р.

Подставим результаты вычислений в формулу и произведем расчет показателя $ROI$:
$$
    ROI = \frac{253516.67 - 143246.61}{143246.61} \cdot 100 \% = 76.98 \%.
$$

Для сопоставления также определим рентабельность инвестиций при условии, что организация-разработчик является резидентом Парка высоких технологий:
$$
    ROI = \frac{371000.00 - 143246.61}{143246.61} \cdot 100 \% = 158.98 \%.
$$

Полученное значение рентабельности инвестиций при обычном режиме налогообложения является существенно положительным, что свидетельствует о заметном превышении годового экономического эффекта над вложениями в разработку и реализацию продукта. Следовательно, разработка и вывод на рынок персонального ассистента организации времени по принятой цене являются экономически целесообразными.

Вариант расчета показателей дисконтированной эффективности инвестиций в данном случае не применяется, поскольку проект окупается менее чем за один год и в соответствии с методикой достаточно определения показателя $ROI$.

\subsection{Вывод по результатам расчета}

В результате выполненного экономического обоснования установлено, что затраты на разработку и реализацию персонального ассистента организации времени составляют 143246.61 р. При принятой отпускной цене лицензии 200 р. в год и прогнозируемом объеме продаж 5300 лицензий годовой прирост чистой прибыли при обычном режиме налогообложения составляет 253516.67 р., а рассчитанная рентабельность инвестиций -- 76.98 \%. Это означает, что проект окупается менее чем за один год и обеспечивает положительный экономический результат для организации-разработчика.

Дополнительный расчет для случая, если организация-разработчик является резидентом Парка высоких технологий, показывает еще более высокий результат: годовой прирост чистой прибыли достигает 371000.00 р., а рентабельность инвестиций -- 158.98 \%. Следовательно, при выбранных исходных предпосылках разработка и реализация программного продукта на массовом рынке являются экономически целесообразными.

К основным факторам, влияющим на полученный результат, относятся выбранная цена лицензии, фактический объем продаж и уровень рентабельности продаж. В связи с этим основными рисками проекта являются более низкий, чем ожидается, спрос на начальном этапе вывода продукта на рынок и необходимость поддержания конкурентоспособности продукта в сегменте цифровых сервисов персональной продуктивности. Вместе с тем выполненные расчеты показывают, что при принятых параметрах коммерциализации проект обладает положительным экономическим потенциалом.

